\chapter{Fonctions caractéristiques et théorème central
limite}\label{ch:b3:clt}

La \omterm{def:b3:probability:space}{loi} des grands nombres dit que les moyennes convergent~; le
théorème central limite dit \emph{comment elles fluctuent}~:
l'erreur, amplifiée par $\sqrt n$, est asymptotiquement
gaussienne --- quelle que soit la \omterm{def:b3:probability:space}{loi} de départ. Cette
universalité est le fait le plus profond des probabilités
élémentaires, et sa preuve naturelle est fourier-analytique~:
la \emph{\omterm{def:b3:clt:cf}{fonction caractéristique}} (la transformée de Fourier
d'une \omterm{def:b3:probability:space}{loi}) convertit les sommes \omterm{def:b3:probability:independence}{indépendantes} en produits, et
la machinerie du~\cref{ch:b3:fouriertransform} ---
injectivité, points fixes \omterm{def:b3:clt:gaussianvector}{gaussiens} --- convertit la
convergence ponctuelle de ces produits en convergence des
\omterm{def:b3:probability:space}{lois} (théorème de Lévy, prouvé en entier). Le chapitre
s'achève sur les \omterm{def:b3:clt:gaussianvector}{vecteurs gaussiens} et la dérivation honnête
des intervalles de confiance utilisés partout en
statistique~; le problème de week-end donne la seconde preuve
de Lindeberg du TCL, avec un taux d'erreur explicite.

\section{Fonctions caractéristiques}

\begin{definition}\label{def:b3:clt:cf}
La \emph{fonction caractéristique}\index{fonction
caractéristique} d'une \omterm{def:b3:probability:space}{variable aléatoire} réelle $X$ est
\[
\varphi_X(\xi) = \E\bigl[\eu^{\iu\xi X}\bigr]
= \int_\R \eu^{\iu\xi x}\,\dd\P_X(x)
\qquad (\xi \in \R)
\]
(le théorème de transfert la calcule depuis la \omterm{def:b3:probability:space}{loi}~; pour une
\omterm{ex:b3:lebesgue:gamma}{densité} $f$, $\varphi_X(\xi) = \hat f(-\xi)$ dans la
convention du~\cref{ch:b3:fouriertransform}).
\end{definition}

\begin{proposition}\label{prop:b3:clt:cfbasics}
(a) $\varphi_X(0) = 1$, $\abs{\varphi_X} \leq 1$, et
$\varphi_X$ est uniformément \omterm{def:b3:topology:continuity}{continue}~;
$\varphi_{aX + b}(\xi) = \eu^{\iu b\xi}\varphi_X(a\xi)$.
(b) Si $X, Y$ sont \emph{\omterm{def:b3:probability:independence}{indépendantes}}~:
$\varphi_{X+Y} = \varphi_X\,\varphi_Y$.
(c) Si $\E\abs X^k < \infty$, alors $\varphi_X \in \mathcal
C^k$ avec $\varphi_X^{(j)}(0) = \iu^j\,\E[X^j]$ pour $j \leq
k$~; en particulier, pour $X$ centrée dans $L^2$ de variance
$\sigma^2$~:
\[
\varphi_X(\xi) = 1 - \frac{\sigma^2\xi^2}{2} +
o(\xi^2) \qquad (\xi \to 0).
\]
(d) Gaussienne~: $X \sim \mathcal N(m, \sigma^2)$ a
$\varphi_X(\xi) = \eu^{\iu m\xi - \sigma^2\xi^2/2}$.
\end{proposition}

\begin{proof}
(a) Les bornes sont immédiates~; continuité~:
$\abs{\varphi(\xi + h) - \varphi(\xi)} \leq
\E\abs{\eu^{\iu hX} - 1} \to 0$ quand $h \to 0$ par
convergence dominée, uniformément en $\xi$. La règle affine
est une substitution. (b) $\eu^{\iu\xi(X+Y)} =
\eu^{\iu\xi X}\eu^{\iu\xi Y}$, et les \omterm{def:b3:probability:space}{espérances} de produits
de variables \omterm{def:b3:probability:independence}{indépendantes} se factorisent
(\cref{thm:b3:probability:independence}, appliqué aux
parties réelles et imaginaires). (c) Dérivation sous
l'\omterm{def:b3:probability:space}{espérance}, dominée par $\E\abs X^j$
(\cref{thm:b3:lebesgue:paramdiff})~; le développement de
Taylor en $0$ est alors Taylor--Young pour la fonction
$\mathcal C^2$ $\varphi$. (d) Pour $\mathcal N(0,1)$~: la
transformée gaussienne
(\cref{ex:b3:fouriertransform:gaussian} avec $a = \frac12$)
donne $\int\eu^{\iu\xi x}\frac{\eu^{-x^2/2}}{\sqrt{2\pi}}\dd
x = \eu^{-\xi^2/2}$~; le cas général par la règle affine.
\end{proof}
\begin{theorem}[Injectivité]\label{thm:b3:clt:injectivity}
Si $\varphi_X = \varphi_Y$, alors $X$ et $Y$ ont la même
\omterm{def:b3:probability:space}{loi}. Plus précisément, pour $N \sim \mathcal N(0,1)$
indépendante de $X$ et $\varepsilon > 0$, la variable lissée
$X + \varepsilon N$ a pour \omterm{ex:b3:lebesgue:gamma}{densité}
\[
p_\varepsilon(x) = \frac1{2\pi}\int_\R
\varphi_X(-\xi)\,\eu^{-\varepsilon^2\xi^2/2}\,
\eu^{\iu\xi x}\,\dd\xi ,
\]
déterminée par $\varphi_X$ seule~; faire $\varepsilon \to 0$
récupère la \omterm{def:b3:probability:space}{loi} de $X$.\index{fonction caractéristique,
injectivité}
\end{theorem}

\begin{proof}
$X + \varepsilon N$ a pour \omterm{ex:b3:lebesgue:gamma}{densité} $p_\varepsilon(x) =
\E\bigl[g_\varepsilon(x - X)\bigr]$, où $g_\varepsilon$ est
la \omterm{ex:b3:lebesgue:gamma}{densité} $\mathcal N(0, \varepsilon^2)$~: en effet pour
$B$ borélien, l'\omterm{def:b3:probability:independence}{indépendance} et Tonelli donnent $\P(X +
\varepsilon N \in B) = \int\!\!\int\mathbf 1_B(x +
\varepsilon n)g_1(n)\,\dd n\,\dd\P_X(x) =
\int_B\E[g_\varepsilon(t - X)]\dd t$ (substituer, puis
Tonelli encore). Écrire $g_\varepsilon$ par inversion de
Fourier de sa transformée
(\cref{exo:b3:fouriertransform:4}, rescalée)~:
$g_\varepsilon(u) = \frac1{2\pi}\int
\eu^{-\varepsilon^2\xi^2/2}\eu^{\iu\xi u}\dd\xi$, et Fubini
(tout est dominé par le facteur \omterm{def:b3:clt:gaussianvector}{gaussien})~:
\[
p_\varepsilon(x) = \frac1{2\pi}\int_\R
\varphi_X(-\xi)\,\eu^{-\varepsilon^2\xi^2/2}\,
\eu^{\iu\xi x}\,\dd\xi ,
\]
une fonctionnelle de $\varphi_X$ seule. Si $\varphi_X =
\varphi_Y$~: $X + \varepsilon N$ et $Y + \varepsilon N$ ont
même \omterm{def:b3:probability:space}{loi} pour tout $\varepsilon$~; pour $f$ \omterm{def:b3:topology:continuity}{continue} bornée,
$\E f(X + \varepsilon N) \to \E f(X)$ quand $\varepsilon \to
0$ (convergence dominée, $X + \varepsilon N \to X$
pointwise sur l'espace produit), donc $\E f(X) = \E f(Y)$
pour toute telle $f$ --- et cela détermine la \omterm{def:b3:probability:space}{loi}~: pour
chaque $t$, encadrer $\mathbf 1_{\intoc{-\infty}t}$ entre
les rampes \omterm{def:b3:topology:continuity}{continues} bornées $f_k^\pm$ (égales à $1$ sur
$\intoc{-\infty}{t \mp \frac1k}$, à $0$ au-delà de $t \pm
\frac1k$, affines entre)~; en passant à la limite dans $\E
f_k^-(X) \leq F_X(t) \leq \E f_k^+(X)$ on obtient $F_X(t)
= F_Y(t)$ en tout $t$ où les deux sont \omterm{def:b3:topology:continuity}{continues}, donc
partout par continuité à droite et \omterm{ex:b3:lebesgue:gamma}{densité} des points de
continuité communs (les deux $F$ ont dénombrablement
beaucoup de sauts)~; des fonctions de répartition égales
forcent des \omterm{def:b3:probability:space}{lois} égales (\cref{exo:b3:measure:3}, reposant
sur le~\cref{thm:b3:measure:uniqueness}).
\end{proof}

\section{Convergence en loi}

\begin{definition}\label{def:b3:clt:cid}
$X_n$ \emph{converge en \omterm{def:b3:probability:space}{loi}} (ou en distribution) vers $X$,
noté $X_n \Rightarrow X$, si\index{convergence en loi}
\[
\E\bigl[f(X_n)\bigr] \longrightarrow \E\bigl[f(X)\bigr]
\qquad\text{pour toute } f\colon\R\to\R \text{ continue
bornée}.
\]
Équivalemment (\cref{exo:b3:clt:4})~: $F_{X_n}(t) \to
F_X(t)$ en tout point de continuité $t$ de $F_X$. Les $X_n$
n'ont pas besoin de vivre sur un \omterm{def:b3:probability:space}{espace probabilisé}
commun~: seules les \omterm{def:b3:probability:space}{lois} comptent.
\end{definition}

\begin{theorem}[Théorème de sélection de Helly]
\label{thm:b3:clt:helly}
Toute suite $(F_n)$ de fonctions de répartition admet une
sous-suite convergeant pointwise, en tout point de continuité
de la limite, vers une $G \colon \R \to \intcc01$ croissante
et \omterm{def:b3:topology:continuity}{continue} à droite --- éventuellement avec $G(+\infty) -
G(-\infty) < 1$ (de la masse peut s'échapper à
l'infini).\index{théorème de Helly}
\end{theorem}

\begin{proof}
L'extraction diagonale donne $F_{n_k}(q) \to \ell(q)$ pour
tout rationnel $q$ (valeurs dans le compact $\intcc01$).
Définir $G(t) = \inf\{\ell(q) : q \in \Q, q > t\}$~:
croissante~; \omterm{def:b3:topology:continuity}{continue} à droite (un infimum sur des \omterm{def:b3:topology:topology}{voisinages}
rationnels rétrécissants depuis la droite). En un point de
continuité $t$ de $G$~: pour des rationnels $q_1 < t < q_2$,
\[
\ell(q_1) \leq \liminf F_{n_k}(t) \leq \limsup F_{n_k}(t)
\leq \ell(q_2),
\]
par monotonie de chaque $F_{n_k}$. Depuis la définition de
$G$ comme infimum et la monotonie de $\ell$ sur les
rationnels~: $G(s) \leq \ell(q) \leq G(q)$ dès que $s < q$.
Prendre $s < q_1 < t$ donne $\ell(q_1) \geq G(s)$, et
$\ell(q_2) \leq G(q_2)$~; en faisant $s \uparrow t$ et $q_2
\downarrow t$, la continuité de $G$ en $t$ force le
$\liminf$ et le $\limsup$ vers $G(t)$.
\end{proof}

\begin{lemma}[Tension depuis la fonction caractéristique]
\label{lem:b3:clt:tightness}
Pour toute \omterm{def:b3:probability:space}{variable aléatoire} $X$ et $u > 0$~:
\[
\P\Bigl(\abs X \geq \frac2u\Bigr) \;\leq\;
\frac1u\int_{-u}^{u}\bigl(1 -
\operatorname{Re}\varphi_X(\xi)\bigr)\,\dd\xi .
\]
\end{lemma}

\begin{proof}
Par Tonelli--Fubini (intégrande bornée, région finie en
$\xi$)~:
\[
\frac1u\int_{-u}^u\bigl(1 -
\operatorname{Re}\varphi_X(\xi)\bigr)\dd\xi
= \E\Bigl[\frac1u\int_{-u}^u(1 - \cos(\xi X))\,\dd\xi\Bigr]
= 2\,\E\Bigl[1 - \frac{\sin(uX)}{uX}\Bigr]
\]
(interpréter le crochet comme sa limite $0$ en $X = 0$).
L'intégrande est positive ($\abs{\sin t} \leq \abs t$), et
pour $\abs{uX} \geq 2$~: $1 - \frac{\sin(uX)}{uX} \geq 1 -
\frac1{\abs{uX}} \geq \frac12$. En ne gardant que
l'événement $\{\abs{uX} \geq 2\}$ dans l'\omterm{def:b3:probability:space}{espérance}, on laisse
au moins $2 \cdot \frac12\,\P(\abs X \geq \frac2u)$, qui
est l'affirmation.
\end{proof}

\begin{theorem}[Théorème de continuité de Lévy]
\label{thm:b3:clt:levy}
Soient $(X_n)$ des \omterm{def:b3:probability:space}{variables aléatoires} dont les \omterm{def:b3:clt:cf}{fonctions
caractéristiques} convergent pointwise~: $\varphi_{X_n}(\xi)
\to \varphi(\xi)$ pour tout $\xi$, où $\varphi = \varphi_X$
est la \omterm{def:b3:clt:cf}{fonction caractéristique} d'une certaine \omterm{def:b3:probability:space}{variable
aléatoire} $X$. Alors $X_n \Rightarrow X$.\index{théorème de
continuité de Lévy}
\end{theorem}

\begin{proof}
\emph{Tension.} Fixer $\varepsilon > 0$. Comme $\varphi$ est
\omterm{def:b3:topology:continuity}{continue} en $0$ avec $\varphi(0) = 1$, choisir $u > 0$ avec
$\frac1u\int_{-u}^u(1 - \operatorname{Re}\varphi) <
\varepsilon$~; par convergence dominée (intégrande bornée
par $2$ sur le $[-u,u]$ fixe), la même intégrale pour
$\varphi_{X_n}$ est $< 2\varepsilon$ pour $n$ grand~: le
\cref{lem:b3:clt:tightness} donne $\P(\abs{X_n} \geq
\frac2u) \leq 2\varepsilon$ pour $n$ grand, et élargir la
constante gère les finiment beaucoup d'autres~: les \omterm{def:b3:probability:space}{lois}
sont \emph{tendues} --- aucune masse ne s'échappe.

\emph{Sous-suites.} Soit $(F_{n_k})$ une sous-suite
quelconque~; par Helly (\cref{thm:b3:clt:helly}) extraire
$F_{n_{k_j}} \to G$ aux points de continuité. La tension
force $G(-\infty) = 0$, $G(+\infty) = 1$ ($G(\frac2u) -
G(-\frac2u) \geq 1 - 2\varepsilon$ aux points de
continuité)~: $G$ est une vraie fonction de répartition,
d'une certaine variable $Y$. Alors $X_{n_{k_j}} \Rightarrow
Y$ (\cref{exo:b3:clt:4}, convergence en \omterm{def:b3:probability:space}{loi} depuis les $F$),
donc $\varphi_{X_{n_{k_j}}} \to \varphi_Y$ \emph{pointwise}
($x \mapsto \eu^{\iu\xi x}$ est \omterm{def:b3:topology:continuity}{continue} bornée, parties
réelle et imaginaire séparément)~; en comparant avec
l'hypothèse~: $\varphi_Y = \varphi = \varphi_X$, et
l'injectivité (\cref{thm:b3:clt:injectivity}) donne $Y \sim
X$, c'est-à-dire $G = F_X$.

\emph{Conclusion.} Toute sous-suite de $(F_n)$ a une
sous-sous-suite convergeant vers la \emph{même} $F_X$ (en
ses points de continuité)~; donc $F_n(t) \to F_X(t)$ en tout
point de continuité $t$ (une suite réelle dont toutes les
sous-suites ont des sous-sous-suites de même limite
converge)~: $X_n \Rightarrow X$.
\end{proof}

\section{Le théorème central limite}

\begin{theorem}[Théorème central limite]\label{thm:b3:clt:clt}
Soient $(X_n)$ i.i.d.\ avec $\E X_1 = m$ et $\V(X_1) =
\sigma^2 \in \intoo0\infty$. Alors\index{théorème central
limite}
\[
\frac{S_n - nm}{\sigma\sqrt n} \;\Longrightarrow\; \mathcal
N(0, 1) :
\qquad
\P\Bigl(a \leq \frac{S_n - nm}{\sigma\sqrt n} \leq
b\Bigr) \longrightarrow
\frac{1}{\sqrt{2\pi}}\int_a^b\eu^{-x^2/2}\,\dd x
\]
pour tous $a < b$.
\end{theorem}

\begin{proof}
Centrer et normaliser~: $Z_i = \frac{X_i - m}{\sigma}$
(i.i.d., moyenne $0$, variance $1$) et $T_n =
\frac1{\sqrt n}\sum_{i\leq n}Z_i$. Par \omterm{def:b3:probability:independence}{indépendance} et la
règle affine (\cref{prop:b3:clt:cfbasics})~:
\[
\varphi_{T_n}(\xi) =
\varphi_{Z}\Bigl(\frac{\xi}{\sqrt n}\Bigr)^{n},
\qquad
\varphi_Z(\eta) = 1 - \frac{\eta^2}2 + \eta^2\rho(\eta),\quad
\rho(\eta)\to0 .
\]
Fixer $\xi$ et poser $a_n = \varphi_Z(\xi/\sqrt n)$, $b_n =
1 - \frac{\xi^2}{2n}$~: tous deux ont un \omterm{def:b3:modules:module}{module} $\leq 1$
pour $n$ grand ($\abs{b_n} \leq 1$ dès que $\xi^2 \leq 4n$~;
$\abs{a_n} \leq 1$ toujours). L'inégalité élémentaire
$\abs{a^n - b^n} \leq n\abs{a - b}$ pour $\abs a, \abs b
\leq 1$ (télescopage $a^n - b^n = \sum a^k(a - b)b^{n-1-k}$)
donne
\[
\Bigl|\varphi_{T_n}(\xi) - \Bigl(1 -
\frac{\xi^2}{2n}\Bigr)^{n}\Bigr|
\leq n\,\Bigl|\varphi_Z\Bigl(\frac\xi{\sqrt n}\Bigr) - 1 +
\frac{\xi^2}{2n}\Bigr|
= \xi^2\,\Bigl|\rho\Bigl(\frac{\xi}{\sqrt n}\Bigr)\Bigr|
\longrightarrow 0,
\]
tandis que $\bigl(1 - \frac{\xi^2}{2n}\bigr)^n \to
\eu^{-\xi^2/2}$ (logarithme réel). Donc $\varphi_{T_n}(\xi)
\to \eu^{-\xi^2/2} = \varphi_{\mathcal N(0,1)}(\xi)$
(\cref{prop:b3:clt:cfbasics}(d)) pour tout $\xi$~: Lévy
(\cref{thm:b3:clt:levy}) conclut $T_n \Rightarrow \mathcal
N(0,1)$. Les probabilités d'intervalles suivent car
$F_{\mathcal N}$ est \omterm{def:b3:topology:continuity}{continue} partout.
\end{proof}

\begin{example}[Intervalles de confiance, honnêtement dérivés]
\label{ex:b3:clt:confidence}
Sonder $n$ électeurs indépendants~; $\hat p_n = S_n/n$
estime le vrai $p$, avec $\sigma^2 = p(1-p) \leq \frac14$.
Le TCL donne, pour $n$ grand,
\[
\P\Bigl(\abs{\hat p_n - p} \leq
\frac{z}{2\sqrt n}\Bigr)
\;\geq\; \P\Bigl(\Bigl|\frac{S_n - np}{\sigma\sqrt n}\Bigr|
\leq z\Bigr)
\longrightarrow \Phi(z) - \Phi(-z),
\]
où $\Phi$ est la fonction de répartition gaussienne
standard. Avec $z = 1.96$~: confiance asymptotique $95\%$,
et une marge $\frac{1.96}{2\sqrt n} \leq 3\%$ exige $n \geq
\bigl(\frac{1.96}{0.06}\bigr)^2 \approx 1068$ --- le nombre
derrière chaque <<~$\pm3$ points, $95\%$~>> qu'on lit~;
comparer aux $5556$ de Bienaymé--Tchebychev
(\cref{exo:b3:probability:7}). Le $\sqrt n$ est
universel~: pour diviser l'erreur par deux, quadrupler
l'échantillon --- la même \omterm{def:b3:probability:space}{loi} qui fixe le coût du
\omterm{ex:b3:probability:sllnapps}{Monte-Carlo} (\cref{exo:b3:clt:7}).
\end{example}

\section{Vecteurs gaussiens}

\begin{definition}\label{def:b3:clt:gaussianvector}
Un vecteur aléatoire $X = (X_1, \dots, X_d)$ est
\emph{gaussien}\index{vecteur gaussien} si toute combinaison
linéaire $\langle t, X\rangle = \sum t_iX_i$ est une
variable gaussienne réelle (éventuellement dégénérée). Sa
\omterm{def:b3:probability:space}{loi} est déterminée par le vecteur moyen $m = (\E X_i)$ et la
\emph{matrice de covariance} $\Sigma =
\bigl(\operatorname{Cov}(X_i, X_j)\bigr)$~: en effet la
\omterm{def:b3:clt:cf}{fonction caractéristique} du vecteur, $\varphi_X(t) =
\E\eu^{\iu\langle t, X\rangle}$, est la valeur en $1$ de la
fc de $\langle t, X\rangle$~:
\[
\varphi_X(t) = \exp\Bigl(\iu\langle t, m\rangle -
\tfrac12\,t^{\mathsf T}\Sigma\,t\Bigr),
\]
et les \omterm{def:b3:clt:cf}{fonctions caractéristiques} $d$-dimensionnelles sont
injectives (même preuve de lissage que
le~\cref{thm:b3:clt:injectivity}, gaussiennes
coordonnée par coordonnée).
\end{definition}

\begin{theorem}\label{thm:b3:clt:gaussianvector}
Soit $X$ un \omterm{def:b3:clt:gaussianvector}{vecteur gaussien}.
\begin{enumerate}
  \item Toute image affine $AX + b$ est un \omterm{def:b3:clt:gaussianvector}{vecteur
        gaussien}.
  \item Les composantes $X_i$ sont \emph{\omterm{def:b3:probability:independence}{indépendantes}} si
        et seulement si $\Sigma$ est diagonale~: pour des
        variables conjointement gaussiennes, non corrélées
        $=$ \omterm{def:b3:probability:independence}{indépendantes}.
  \item Si $\Sigma$ est inversible, $X$ a pour \omterm{ex:b3:lebesgue:gamma}{densité}
        $\frac{1}{(2\pi)^{d/2}\sqrt{\det\Sigma}}
        \exp\bigl(-\frac12(x - m)^{\mathsf T}\Sigma^{-1}(x -
        m)\bigr)$.
\end{enumerate}
\end{theorem}

\begin{proof}
(1) Les combinaisons linéaires des composantes de $AX + b$
sont des fonctions affines de combinaisons linéaires de
$X$~: gaussiennes (l'image affine d'une variable gaussienne
est gaussienne).
(2) Si $\Sigma$ est diagonale, la \omterm{def:b3:clt:cf}{fonction caractéristique}
se factorise~: $\varphi_X(t) = \prod_i\exp(\iu t_im_i -
\frac12\Sigma_{ii}t_i^2) = \prod\varphi_{X_i}(t_i)$, qui est
la \omterm{def:b3:clt:cf}{fonction caractéristique} de la \omterm{def:b3:probability:space}{loi} produit
(\cref{thm:b3:probability:independence} lue via
l'injectivité $d$-dimensionnelle)~: les composantes sont
\omterm{def:b3:probability:independence}{indépendantes}. La réciproque est l'annulation des
covariances de variables $L^2$ \omterm{def:b3:probability:independence}{indépendantes}.
(3) Diagonaliser $\Sigma = P D P^{\mathsf T}$ ($P$
orthogonal, $D > 0$ diagonale ---
\cref{exo:b3:submanifolds:8})~; le vecteur $Y = P^{\mathsf
T}(X - m)$ est \omterm{def:b3:clt:gaussianvector}{gaussien} de covariance $D$~: par (2) ses
composantes sont \omterm{def:b3:probability:independence}{indépendantes} $\mathcal N(0, d_i)$, donc
$Y$ a la \omterm{ex:b3:lebesgue:gamma}{densité} produit~; pousser en avant par le
$x = m + PY$ préservant le volume
(\cref{thm:b3:product:linearchange}, $\abs{\det P} = 1$) et
réécrire l'exposant de façon invariante.
\end{proof}

\begin{theorem}[TCL multidimensionnel]
\label{thm:b3:clt:multiclt}
Soient $(X_n)$ des \emph{vecteurs} aléatoires i.i.d.\ de
$\R^d$ de carré \omterm{def:b3:lebesgue:l1}{intégrable}, d'\omterm{def:b3:probability:space}{espérance} $m$ et de matrice
de covariance $\Sigma$. Alors $\frac{S_n - nm}{\sqrt n}$
converge en \omterm{def:b3:probability:space}{loi} vers le \omterm{def:b3:clt:gaussianvector}{vecteur gaussien} $\mathcal N(0,
\Sigma)$.\index{théorème central limite!multidimensionnel}
\admitted
\end{theorem}

\begin{remark}\label{rem:b3:clt:multiclt}
Presque tout est déjà entre nos mains. Pour chaque
direction $t \in \R^d$, la variable réelle $\langle t,
\frac{S_n - nm}{\sqrt n}\rangle$ est une somme normalisée
de variables réelles i.i.d.\ de variance $t^{\mathsf
T}\Sigma t$, donc le calcul du~\cref{thm:b3:clt:clt} donne
la convergence ponctuelle des \omterm{def:b3:clt:cf}{fonctions caractéristiques}
$d$-dimensionnelles vers $\eu^{-t^{\mathsf T}\Sigma t/2}$,
la \omterm{def:b3:clt:cf}{fonction caractéristique} de $\mathcal N(0, \Sigma)$
(\cref{def:b3:clt:gaussianvector}). Ce que nous n'avons pas
reprouvé est le théorème de continuité de Lévy \emph{dans
$\R^d$}~: la sélection de Helly et l'estimée de tension se
généralisent routinièrement (coordonnée par coordonnée), et
cette réduction de Cramér--Wold est menée honnêtement dans
tout cours de probabilités de master~; rien au-delà des
méthodes de ce chapitre n'est nécessaire.
\end{remark}

\begin{method}\label{met:b3:clt:method}
Pour identifier une \omterm{def:b3:probability:space}{loi} limite~: calculer les \omterm{def:b3:clt:cf}{fonctions
caractéristiques}, prendre la limite ponctuelle, la
reconnaître (gaussienne $\eu^{-\sigma^2\xi^2/2}$, Poisson
$\eu^{\lambda(\eu^{\iu\xi}-1)}$, exponentielle
$\frac{\lambda}{\lambda - \iu\xi}$, $\dots$) et invoquer
Lévy. Le rituel en trois étapes (\omterm{def:b3:probability:independence}{indépendance} $\to$
produit~; Taylor en $0$ $\to$ limite exponentielle~; Lévy
$\to$ convergence en \omterm{def:b3:probability:space}{loi}) prouve le TCL, la \omterm{def:b3:probability:space}{loi} de Poisson
des événements rares (\cref{exo:b3:clt:5}), et tout
théorème limite classique de ce cours. Pour les énoncés
p.s., revenir à la boîte à outils du
\cref{ch:b3:probability}~: les deux chapitres répondent à
des questions différentes sur le même $S_n$.
\end{method}


\section{Exercices}

\begin{exercise}[$\star$]\label{exo:b3:clt:1}
Calculer les \omterm{def:b3:clt:cf}{fonctions caractéristiques}~: uniforme sur
$\intcc{-1}1$~; exponentielle $\mathcal E(\lambda)$~;
Poisson $\mathcal P(\lambda)$~; binomiale $\mathcal B(n,
p)$. En déduire via le~\cref{thm:b3:clt:injectivity} que la
somme de variables de Poisson \omterm{def:b3:probability:independence}{indépendantes} ($\lambda,
\mu$) est Poisson $(\lambda + \mu)$.
\end{exercise}

\begin{exercise}[$\star\star$]\label{exo:b3:clt:2}
(a) Montrer que $\varphi_X$ est à valeurs réelles si et
seulement si $X$ et $-X$ ont la même \omterm{def:b3:probability:space}{loi} (une variable
\emph{symétrique}).
(b) Supposer $\abs{\varphi_X(\xi_0)} = 1$ pour un $\xi_0
\neq 0$. Montrer que $X$ est presque sûrement supportée sur
une progression arithmétique $a + \frac{2\pi}{\xi_0}\Z$
\emph{(écrire $\varphi_X(\xi_0) = \eu^{\iu\theta}$ et
calculer $\E[1 - \cos(\xi_0X - \theta)]$)}. En déduire que
si $X$ a une \omterm{ex:b3:lebesgue:gamma}{densité}, alors $\abs{\varphi_X(\xi)} < 1$ pour
tout $\xi \neq 0$.
\end{exercise}

\begin{exercise}[$\star\star$]\label{exo:b3:clt:3}
Soient $X \sim \mathcal N(m_1, \sigma_1^2)$ et $Y \sim
\mathcal N(m_2, \sigma_2^2)$ \omterm{def:b3:probability:independence}{indépendantes}. Montrer $X + Y
\sim \mathcal N(m_1 + m_2, \sigma_1^2 + \sigma_2^2)$, et
plus généralement que la famille gaussienne est \omterm{def:b3:ode:stability}{stable} sous
sommes \omterm{def:b3:probability:independence}{indépendantes} et applications affines. Contraste~: la
somme de deux gaussiennes \emph{dépendantes} est-elle
toujours gaussienne~? (\cref{exo:b3:clt:9}.)
\end{exercise}

\begin{exercise}[$\star\star$]\label{exo:b3:clt:4}
(a) Prouver l'équivalence dans
la~\cref{def:b3:clt:cid}~: si $\E f(X_n) \to \E f(X)$ pour
toute $f$ \omterm{def:b3:topology:continuity}{continue} bornée, alors $F_{X_n}(t) \to F_X(t)$
aux points de continuité \emph{(encadrer $\mathbf
1_{\intoc{-\infty}t}$ entre deux rampes en escalier
\omterm{def:b3:topology:continuity}{continues})}~; et réciproquement \emph{(approximer une $f$
\omterm{def:b3:topology:continuity}{continue} bornée par des sommes de fonctions en rampe, ou
conditionner sur une grille fine de points de continuité)}
--- la réciproque peut être traitée d'abord pour $f$
uniformément \omterm{def:b3:topology:continuity}{continue}, puis en général.
(b) Montrer que $X_n \Rightarrow c$ (une constante) implique
$X_n \to c$ en probabilité.
\end{exercise}

\begin{exercise}[$\star\star$]\label{exo:b3:clt:5}
(\omterm{def:b3:probability:space}{Loi} des événements rares) Soit $X_n \sim \mathcal B(n,
p_n)$ avec $np_n \to \lambda > 0$. Montrer, via les
\omterm{def:b3:clt:cf}{fonctions caractéristiques} et le~\cref{thm:b3:clt:levy},
que $X_n \Rightarrow \mathcal P(\lambda)$. Vérification
numérique~: comparer $\P(X = 0)$ pour $\mathcal B(100,
0.02)$ et $\mathcal P(2)$.
\end{exercise}

\begin{exercise}[$\star\star$]\label{exo:b3:clt:6}
(a) Un dé équitable est lancé $n = 1000$ fois~; approximer
la probabilité que le total dépasse $3600$ (moyenne $3500$,
variance par lancer $\frac{35}{12}$).
(b) Pour $S \sim \mathcal B(100, \frac12)$, approximer
$\P(45 \leq S \leq 55)$ par le TCL avec la correction de
continuité ($\pm\frac12$), et commenter l'effet de la
correction.
\end{exercise}

\begin{exercise}[$\star\star$]\label{exo:b3:clt:7}
(Erreur de \omterm{ex:b3:probability:sllnapps}{Monte-Carlo}) Dans le cadre du
\cref{pb:b3:probability:1}, question~11, avec $g \in
L^2(\intcc01^d)$, soit $\sigma^2 = \V(g(U_1))$ et $I =
\int g$. Montrer
\[
\sqrt n\,\Bigl(\frac1n\sum_{k\leq n}g(U_k) - I\Bigr)
\Longrightarrow \mathcal N(0, \sigma^2),
\]
et en déduire la barre d'erreur asymptotique $95\%$ $\pm
1.96\,\sigma/\sqrt n$ --- indépendante de la dimension $d$.
Comparer avec la règle du point milieu déterministe en
dimension $d$ (erreur $\sim n^{-2/d}$ pour des intégrandes
$\mathcal C^2$)~: à partir de quelle dimension
l'échantillonnage aléatoire l'emporte-t-il~?
\end{exercise}

\begin{exercise}[$\star\star\star$]\label{exo:b3:clt:8}
(Slutsky) Supposer $X_n \Rightarrow X$ et $Y_n \to c$ en
probabilité ($c$ constante). Montrer $X_n + Y_n \Rightarrow
X + c$ et $Y_nX_n \Rightarrow cX$. \emph{(Travailler avec
les \omterm{def:b3:clt:cf}{fonctions caractéristiques} et la borne
$\abs{\E\eu^{\iu\xi(X_n+Y_n)} - \eu^{\iu\xi
c}\E\eu^{\iu\xi X_n}} \leq \E\abs{\eu^{\iu\xi(Y_n - c)} -
1}$, scinder sur $\abs{Y_n - c} \leq \delta$.)}
Application~: dans l'\cref{ex:b3:clt:confidence}, justifier
le remplacement de l'inconnue $\sigma = \sqrt{p(1-p)}$ par
$\sqrt{\hat p_n(1 - \hat p_n)}$.
\end{exercise}

\begin{exercise}[$\star\star\star$]\label{exo:b3:clt:9}
Soit $X \sim \mathcal N(0,1)$ et $\varepsilon$ indépendante
avec $\P(\varepsilon = \pm1) = \frac12$~; poser $Y =
\varepsilon X$.
(a) Montrer $Y \sim \mathcal N(0,1)$ et
$\operatorname{Cov}(X, Y) = 0$.
(b) Montrer que $X$ et $Y$ ne sont \emph{pas}
\omterm{def:b3:probability:independence}{indépendantes}, et que $(X, Y)$ n'est pas un \omterm{def:b3:clt:gaussianvector}{vecteur
gaussien} \emph{(calculer $\P(X + Y = 0)$)}.
(c) Morale~: le~\cref{thm:b3:clt:gaussianvector}(2) exige
la gaussianité conjointe --- <<~gaussiennes non corrélées~>>
seules ne prouvent rien.
\end{exercise}

\begin{exercise}[$\star\star$]\label{exo:b3:clt:10}
La \omterm{def:b3:probability:space}{loi} de Cauchy a pour \omterm{ex:b3:lebesgue:gamma}{densité} $\frac1{\pi(1 + x^2)}$.
(a) Montrer que sa \omterm{def:b3:clt:cf}{fonction caractéristique} est
$\eu^{-\abs\xi}$ (\cref{exo:b3:fouriertransform:1} et
inversion).
(b) Montrer que si $X_1, \dots, X_n$ sont i.i.d.\ Cauchy,
alors $\frac{S_n}n$ est encore Cauchy --- la \emph{même}
\omterm{def:b3:probability:space}{loi}~: la moyenne ne se concentre jamais.
(c) Réconcilier avec les \omterm{def:b3:probability:space}{lois} des grands nombres et le
TCL~: quelles hypothèses échouent~? (Calculer $\E\abs{X_1}$.)
\end{exercise}

\begin{exercise}[$\star\star$]\label{exo:b3:clt:11}
(\omterm{def:b3:probability:space}{Lois} \omterm{def:b3:ode:stability}{stables} en germe) Soient $(X_n)$ i.i.d.\ Cauchy
standard (\cref{exo:b3:clt:10}).
(a) Montrer que pour tous $a, b > 0$, $aX_1 + bX_2$ a la
\omterm{def:b3:probability:space}{loi} de $(a + b)X_1$~: la famille de Cauchy est
\emph{strictement \omterm{def:b3:ode:stability}{stable}} d'\omterm{def:b3:holomorphic:index}{indice} $1$.
(b) Montrer que la famille gaussienne est strictement
\omterm{def:b3:ode:stability}{stable} d'\omterm{def:b3:holomorphic:index}{indice} $2$~: $aX_1 + bX_2 \sim \sqrt{a^2 +
b^2}\,X_1$ pour $X_i$ i.i.d.\ $\mathcal N(0,1)$.
(c) Expliquer, via les \omterm{def:b3:clt:cf}{fonctions caractéristiques} de la
forme $\eu^{-c\abs\xi^\alpha}$, pourquoi la stabilité
d'\omterm{def:b3:holomorphic:index}{indice} $\alpha$ force la normalisation $n^{1/\alpha}$
pour les sommes, et ce que cela dit sur les bassins
d'attraction du TCL~: quelles sommes i.i.d.\ peuvent
converger, après normalisation affine, vers une \omterm{def:b3:probability:space}{loi} de
Cauchy plutôt qu'une gaussienne~?
\end{exercise}

\begin{exercise}[$\star\star$]\label{exo:b3:clt:12}
(La fonction de répartition empirique) Soient $(X_n)$
i.i.d.\ de fonction de répartition $F$, et $F_n(t) =
\frac1n\#\{k \leq n : X_k \leq t\}$.
(a) Fixer $t$. Montrer que $n F_n(t) \sim \mathcal B(n,
F(t))$, que $F_n(t) \to F(t)$ p.s.\
(\cref{thm:b3:probability:slln}), et que
\[
\sqrt n\,\bigl(F_n(t) - F(t)\bigr) \Longrightarrow
\mathcal N\bigl(0,\ F(t)(1 - F(t))\bigr) .
\]
(b) En quel $t$ la variance asymptotique est-elle
maximale~? Interpréter~: la médiane est l'endroit où une
répartition empirique est le plus difficile à épingler.
(c) Pour $F$ \omterm{def:b3:topology:continuity}{continue}, montrer que la \omterm{def:b3:probability:space}{loi} de
$\sup_t\abs{F_n(t) - F(t)}$ ne dépend pas de $F$
\emph{(se ramener à des variables uniformes via
l'\cref{exo:b3:probability:1})} --- le miracle sans
distribution derrière le test de Kolmogorov--Smirnov~; aucun
calcul de cette \omterm{def:b3:probability:space}{loi} n'est demandé.
\end{exercise}

\section{Problème~: la preuve de Lindeberg du TCL, avec un
taux}

\begin{problem}[{Problème de week-end --- la méthode de
remplacement}]\label{pb:b3:clt:1}
Lindeberg (1922) a prouvé le théorème central limite par une
idée d'une simplicité désarmante~: \emph{échanger les
summands un par un contre des gaussiennes} et contrôler
chaque échange par un développement de Taylor. La méthode
n'a besoin d'aucune analyse de Fourier, produit un taux
d'erreur explicite, et alimente aujourd'hui les preuves
d'universalité à travers la théorie des probabilités.
Soient $(X_i)$ i.i.d., centrées, $\V(X_1) = 1$, avec $\beta
= \E\abs{X_1}^3 < \infty$~; soient $(N_i)$ i.i.d.\
$\mathcal N(0,1)$, \omterm{def:b3:probability:independence}{indépendantes} des $X_i$ (existence~:
\cref{thm:b3:probability:existence}). Poser
\[
T_n = \frac{X_1 + \dots + X_n}{\sqrt n},
\qquad
G_n = \frac{N_1 + \dots + N_n}{\sqrt n} \sim \mathcal N(0,1).
\]

\medskip
\noindent\textbf{Partie I --- L'identité d'échange.} Fixer
$f \in \mathcal C^3_b(\R)$ (trois dérivées \omterm{def:b3:topology:continuity}{continues}
bornées~; $M_3 = \sup\abs{f'''}$). Pour $0 \leq i \leq n$
définir les sommes hybrides
\[
H_i = \frac{X_1 + \dots + X_i + N_{i+1} + \dots +
N_n}{\sqrt n},
\]
donc $H_n = T_n$ et $H_0 = G_n$.
\begin{enumerate}
  \item Écrire $H_i = W_i + \frac{X_i}{\sqrt n}$ et
        $H_{i-1} = W_i + \frac{N_i}{\sqrt n}$ avec $W_i =
        \frac{1}{\sqrt n}\bigl(\sum_{j<i}X_j +
        \sum_{j>i}N_j\bigr)$, et noter que $W_i$ est
        indépendante de la paire $(X_i, N_i)$. Justifier.
  \item Taylor avec reste intégral ou de Lagrange~: pour
        tous réels $w, h$~:
        \[
        \Bigl|f(w + h) - f(w) - f'(w)h -
        \tfrac12f''(w)h^2\Bigr| \leq
        \frac{M_3\,\abs h^3}{6} .
        \]
  \item Appliquer la question~2 deux fois ($h =
        \frac{X_i}{\sqrt n}$ et $h = \frac{N_i}{\sqrt n}$ en
        $w = W_i$), prendre les \omterm{def:b3:probability:space}{espérances}, et utiliser
        l'\omterm{def:b3:probability:independence}{indépendance} plus l'égalité des deux premiers
        moments de $X_i$ et $N_i$ pour montrer
        \[
        \bigl|\E f(H_i) - \E f(H_{i-1})\bigr|
        \leq \frac{M_3}{6}\cdot
        \frac{\beta + \gamma}{n^{3/2}},
        \qquad \gamma = \E\abs{N_1}^3 =
        \frac{2\sqrt2}{\sqrt\pi} .
        \]
  \item Télescoper sur $i$ et conclure la \emph{borne de
        Lindeberg}~:
        \[
        \bigl|\E f(T_n) - \E f(G_n)\bigr| \leq
        \frac{M_3\,(\beta + \gamma)}{6\,\sqrt n} .
        \]
\end{enumerate}

\medskip
\noindent\textbf{Partie II --- Des $f$ lisses au TCL.}
\begin{enumerate}[resume]
  \item Montrer que $\E f(T_n) \to \E f(N)$ pour toute $f
        \in \mathcal C_b^3$, et monter à toutes les $f$
        \omterm{def:b3:topology:continuity}{continues} bornées~: étant donnés une telle $f$ et
        $\varepsilon$, construire $f_\varepsilon \in
        \mathcal C^3_b$ avec $\norm{f - f_\varepsilon}_\infty
        \leq \varepsilon$ sur un grand intervalle --- p.ex.\
        convoler $f$ avec une bosse $\mathcal C^\infty$
        (\cref{thm:b3:lp:regularization}) --- et traiter les
        queues par la tension ($\V(T_n) = 1$ et
        Bienaymé--Tchebychev). Conclure $T_n \Rightarrow
        \mathcal N(0, 1)$~: le théorème central limite,
        reprouvé.
  \item Où la preuve a-t-elle utilisé que les $X_i$ sont
        \emph{identiquement} distribuées~? Montrer que
        c'est à peine le cas~: énoncer et prouver la
        version pour des $X_i$ \omterm{def:b3:probability:independence}{indépendantes}, centrées, non
        identiques avec $\sum_i\V(X_i) = s_n^2$ et moments
        d'ordre trois, obtenant l'erreur
        $\frac{M_3}{6s_n^3}\sum_i\bigl(\E\abs{X_i}^3 +
        \V(X_i)^{3/2}\gamma\bigr)$ --- le vrai théorème de
        Lindeberg sous sa forme de Lyapunov.
\end{enumerate}

\medskip
\noindent\textbf{Partie III --- Dividendes quantitatifs.}
\begin{enumerate}[resume]
  \item (Fonctions de répartition) Soit $t \in \R$ et
        approximer $\mathbf 1_{\intoc{-\infty}t}$ par
        dessus et par dessous par des rampes $\mathcal
        C^3_b$ de largeur $\delta$ (les construire, avec
        $M_3 = O(\delta^{-3})$). En combinant avec la
        Partie~I, dériver la borne à deux termes
        \[
        \sup_{t\in\R}\,\bigl|\P(T_n \leq t) -
        \Phi(t)\bigr| \;\leq\;
        \frac{C_1(\beta + \gamma)}{\delta^3\sqrt n} +
        C_2\,\delta
        \qquad (\text{tout } \delta > 0),
        \]
        avec des constantes explicites (le terme $C_2\delta$
        utilise que $\Phi$ a une \omterm{ex:b3:lebesgue:gamma}{densité} bornée par
        $\frac1{\sqrt{2\pi}}$), et optimiser $\delta \sim
        n^{-1/8}$ pour obtenir un taux uniforme d'ordre
        $n^{-1/8}$.
        (L'optimal $n^{-1/2}$ --- Berry--Esseen --- demande
        des outils plus fins~; le point est un taux
        \emph{explicite} par échange élémentaire.)
  \item (De Moivre--Laplace, quantifié) Spécialiser à $X_i
        = 2B_i - 1$ (signes de pièces équitables)~: comparer
        la conclusion avec l'estimée locale du
        \cref{pb:b3:product:1}, question~7 --- que donne
        chaque méthode que l'autre ne donne pas~?
  \item (Universalité) Expliquer en un paragraphe pourquoi
        la méthode de remplacement montre plus que le
        TCL~: toute statistique de la forme $\E
        f(\text{somme})$ avec $f$ lisse est insensible, à
        l'ordre $n^{-1/2}$, à la \emph{\omterm{def:b3:probability:space}{loi} entière} des
        summands au-delà de ses deux premiers moments ---
        le «~principe d'invariance~» qui sous-tend les
        résultats modernes d'universalité (matrices
        aléatoires, polynômes aléatoires), dont le TCL est
        la première instance.
\end{enumerate}

\medskip
\noindent\textbf{Partie IV --- Lissage poussé~: meilleurs
taux.} La perte de $n^{-1/2}$ ($f$ lisse) à $n^{-1/8}$
(fonctions de répartition) venait de facturer $f'''$ en
norme sup. Les hybrides peuvent en réparer une partie~: ils
contiennent des summands \omterm{def:b3:clt:gaussianvector}{gaussiens}, et les gaussiennes
\emph{lissent}.
\begin{enumerate}[resume]
  \item (Une gaussienne cachée) Pour $1 \leq i \leq n - 1$,
        $h = \frac{X_i}{\sqrt n}$ ou $\frac{N_i}{\sqrt n}$,
        et $\theta \in \intcc01$, écrire $W_i + \theta h =
        A + Z$ avec $Z = \frac{N_{i+1} + \dots + N_n}{\sqrt
        n}$. Montrer que $Z \sim \mathcal N\bigl(0,
        \frac{n-i}n\bigr)$ est indépendante de la paire
        $(A, h)$, et en déduire, pour toute $g$ \omterm{def:b3:topology:continuity}{continue}
        dans $L^1(\R)$,
        \[
        \E\bigl[\abs h^3\,\abs{g(W_i + \theta h)}\bigr]
        \;\leq\; \sqrt{\frac{n}{2\pi(n - i)}}\;
        \norm{g}_{L^1}\;\E\abs h^3 .
        \]
  \item Combiner la question~10 avec la forme intégrale du
        reste de Taylor,
        \[
        f(w + h) = f(w) + f'(w)h + \tfrac12f''(w)h^2 +
        \int_0^1\frac{(1 - \theta)^2}2\,f'''(w + \theta
        h)\,h^3\,\dd\theta,
        \]
        pour refaire les questions~3--4~: pour $f \in
        \mathcal C^3_b$ avec de plus $f''' \in L^1(\R)$,
        \[
        \bigl|\E f(T_n) - \E f(G_n)\bigr| \leq
        \frac{\beta + \gamma}{3\sqrt{2\pi}}\cdot
        \frac{\norm{f'''}_{L^1}}{\sqrt n} +
        \frac{M_3(\beta + \gamma)}{6\,n^{3/2}}
        \]
        \emph{(la question~10 traite les échanges $i \leq n
        - 1$ --- utiliser $\sum_{m=1}^{n-1}m^{-1/2} \leq
        2\sqrt n$ --- et la borne grossière de la question~3
        traite le dernier)}. Vérifier que les rampes de la
        question~7 satisfont $\norm{\psi_\delta'''}_{L^1} =
        K_1\delta^{-2}$ tandis que $M_3 = K\delta^{-3}$,
        les injecter, et optimiser $\delta$~: le taux
        uniforme de fonction de répartition s'améliore en
        $O(n^{-1/6})$.
  \item (Apparier un moment de plus) Supposer de plus $\E
        X_1^3 = 0$ et $\beta_4 = \E X_1^4 < \infty$.
        Calculer $\E N_1^3$ et $\E N_1^4$, développer à
        l'ordre quatre, et prouver le long des mêmes lignes
        que le taux de fonction de répartition devient
        $O(n^{-1/4})$ \emph{(maintenant
        $\norm{\psi_\delta^{(4)}}_{L^1} = K_2\delta^{-3}$
        et $M_4 = K'\delta^{-4}$~; choisir $\delta =
        n^{-1/4}$)}.
  \item (L'obstruction) Supposer que les $k$ premiers
        moments de $X_1$ coïncident avec les \omterm{def:b3:clt:gaussianvector}{gaussiens} ($k
        = 2$ toujours~; $k = 3$ exactement quand $\E X_1^3
        = 0$~; $k \geq 4$ essentiellement jamais, car $\E
        N_1^4 = 3$). Vérifier que le schéma des
        questions~10--12 délivre le taux de fonction de
        répartition $n^{-(k-1)/(2k+2)}$, en équilibrant
        $\delta^{-k}n^{-(k-1)/2}$ contre $\delta$, et
        observer que l'exposant n'approche la valeur de
        Berry--Esseen $\frac12$ que quand $k \to \infty$.
        Expliquer en quelques phrases pourquoi la méthode
        d'échange sature~: chaque échange est facturé en
        valeur absolue, tandis que la voie de Fourier
        (inégalité de lissage d'Esseen) exploite
        l'oscillation de la différence des \omterm{def:b3:clt:cf}{fonctions
        caractéristiques} et atteint $C\beta n^{-1/2}$ avec
        seulement trois moments.
\end{enumerate}

\medskip
\noindent\textbf{Partie V --- Deux dimensions~: le TCL
multidimensionnel, par échange.} Soient maintenant les
$X_i$ des \emph{vecteurs} aléatoires i.i.d.\ centrés de
$\R^2$ de matrice de covariance $\Sigma$ et $\beta' =
\E\norm{X_1}^3 < \infty$ (norme euclidienne).
\begin{enumerate}[resume]
  \item (\omterm{def:b3:clt:gaussianvector}{Vecteurs gaussiens}, en ordre) Diagonaliser $\Sigma
        = PDP^{\mathsf T}$ (\cref{exo:b3:submanifolds:8}) et
        poser $C = P\sqrt DP^{\mathsf T}$. Pour $Z = (Z^1,
        Z^2)$ une paire de gaussiennes standard
        \omterm{def:b3:probability:independence}{indépendantes} (\cref{thm:b3:probability:existence}),
        montrer que $N = CZ$ est un \omterm{def:b3:clt:gaussianvector}{vecteur gaussien}
        (\cref{def:b3:clt:gaussianvector}) d'\omterm{def:b3:probability:space}{espérance} $0$,
        de covariance $\Sigma$, avec $\gamma' = \E\norm N^3
        < \infty$~; et que $G_n = \frac{N_1 + \dots +
        N_n}{\sqrt n}$ a pour \omterm{def:b3:probability:space}{loi} $\mathcal N(0, \Sigma)$
        \emph{exactement} pour des copies i.i.d.\ $N_i$.
  \item (Taylor en deux variables) Pour $f \colon \R^2 \to
        \R$ de classe $\mathcal C^3$ avec $M_3 =
        \max_{\abs\alpha = 3}\sup\abs{\partial^\alpha f} <
        \infty$, prouver
        \[
        \Bigl|f(w + h) - f(w) - \langle\nabla f(w),
        h\rangle - \tfrac12\langle h, D^2f(w)\,h\rangle
        \Bigr| \leq \frac{M_3}6\,\bigl(\abs{h_1} +
        \abs{h_2}\bigr)^3 \leq \frac{\sqrt2\,M_3}3\,
        \norm h^3
        \]
        \emph{(étudier $t \mapsto f(w + th)$ sur
        $\intcc01$)}.
  \item (Le TCL dans $\R^2$) Faire tourner le schéma de
        remplacement sur les hybrides vectoriels $H_i$~:
        montrer que les termes d'ordre un et deux
        s'annulent (moyennes et covariances coïncident),
        télescoper, et monter comme à la question~5 (tension
        depuis $\E\norm{T_n}^2 = \operatorname{tr}\Sigma$~;
        mollification maintenant dans $\R^2$,
        \cref{thm:b3:lp:regularization}) pour conclure~:
        pour toute $f \colon \R^2 \to \R$ \omterm{def:b3:topology:continuity}{continue} bornée,
        \[
        \E\,f\Bigl(\frac{X_1 + \dots + X_n}{\sqrt n}\Bigr)
        \longrightarrow \E\,f(N), \qquad N \sim \mathcal
        N(0, \Sigma) :
        \]
        le~\cref{thm:b3:clt:multiclt} en dimension~$2$, avec
        un taux pour les $f$ lisses et sans analyse de
        Fourier.
  \item (Cramér--Wold, et une fluctuation jointe) En
        déduire que $\langle t, \frac{S_n}{\sqrt n}\rangle
        \Rightarrow \mathcal N(0, t^{\mathsf T}\Sigma t)$
        pour tout $t \in \R^2$ fixé. Application~: pour des
        $(\xi_i)$ réels i.i.d., centrés, $\E\xi_1^2 = 1$,
        $\E\xi_1^6 < \infty$ (de sorte que la Partie~V
        s'applique à $V_i = (\xi_i, \xi_i^2 - 1)$), montrer
        \[
        \frac1{\sqrt n}\Bigl(\sum_{i\leq n}\xi_i,\
        \sum_{i\leq n}(\xi_i^2 - 1)\Bigr) \Longrightarrow
        \mathcal N\Bigl(0, \begin{pmatrix} 1 & \E\xi_1^3\\
        \E\xi_1^3 & \E\xi_1^4 - 1\end{pmatrix}\Bigr) :
        \]
        moyenne empirique et second moment empirique
        fluctuent conjointement de façon gaussienne ---
        indépendamment à la limite si et seulement si
        $\E\xi_1^3 = 0$ (\cref{thm:b3:clt:gaussianvector}).
\end{enumerate}

\medskip
\noindent\textbf{Partie VI --- La méthode delta.}
\begin{enumerate}[resume]
  \item Soient $(\hat\theta_n)$ des \omterm{def:b3:probability:space}{variables aléatoires}
        avec $\sqrt n(\hat\theta_n - \theta) \Rightarrow
        \mathcal N(0, \sigma^2)$ pour un paramètre réel
        $\theta$, et $g$ différentiable en $\theta$.
        Prouver la \emph{méthode delta}~:
        \[
        \sqrt n\bigl(g(\hat\theta_n) - g(\theta)\bigr)
        \Longrightarrow \mathcal N\bigl(0,
        g'(\theta)^2\sigma^2\bigr)
        \]
        \emph{(écrire $g(x) - g(\theta) = (g'(\theta) +
        \eta(x))(x - \theta)$ avec $\eta \to 0$ en
        $\theta$~; montrer $\hat\theta_n \to \theta$, puis
        $\eta(\hat\theta_n) \to 0$, en probabilité~;
        conclure avec Slutsky, \cref{exo:b3:clt:8}, et
        \cref{exo:b3:clt:4}(b))}.
  \item Applications. (a) Pour des $(\xi_i)$ réels i.i.d.\
        d'\omterm{def:b3:probability:space}{espérance} $\mu$ et de variance $\sigma^2$, et
        $\bar X_n = \frac1n\sum_{i\leq n}\xi_i$~: montrer
        $\sqrt n(\bar X_n^2 - \mu^2) \Rightarrow \mathcal
        N(0, 4\mu^2\sigma^2)$ quand $\mu \neq 0$, et que
        pour $\mu = 0$ l'énoncé correct vit à une autre
        échelle~: $n\bar X_n^2 \Rightarrow \sigma^2N^2$
        avec $N \sim \mathcal N(0,1)$ (identifier la
        fonction de répartition de la limite). (b)
        (Stabilisation de la variance) Pour $\hat p_n$ la
        fréquence de succès d'un échantillon $\mathcal
        B(1, p)$, $p \in \intoo01$~: montrer que $g(p) =
        \arcsin\sqrt p$ satisfait
        \[
        \sqrt n\,\bigl(g(\hat p_n) - g(p)\bigr)
        \Longrightarrow \mathcal N\Bigl(0, \frac14\Bigr)
        \]
        \emph{quel que soit} $p$ --- une barre d'erreur
        asymptotique \omterm{def:b3:modules:free}{libre} du paramètre inconnu~; comparer
        avec l'\cref{ex:b3:clt:confidence}.
\end{enumerate}

\medskip
\noindent\textbf{Partie VII --- Poisson, par la même
méthode~: le théorème de Le Cam.} Le remplacement connaît
une seconde classe d'universalité~: les sommes de beaucoup
d'événements indépendants \emph{rares}. Pour les \omterm{def:b3:probability:space}{lois} sur
$\N$ la bonne distance est la \emph{variation totale},
\[
d_{\mathrm{TV}}(\mu, \nu) = \sup_{A\subseteq\N}\,
\abs{\mu(A) - \nu(A)} .
\]
\begin{enumerate}[resume]
  \item Montrer que $d_{\mathrm{TV}}(\mu, \nu) =
        \frac12\sum_{k\geq0}\abs{\mu(\{k\}) -
        \nu(\{k\})}$, et prouver la borne de couplage~: pour
        \emph{toute} paire $(X, Y)$ de \omterm{def:b3:probability:space}{variables aléatoires}
        de \omterm{def:b3:probability:space}{lois} $\mu$ et $\nu$ sur le même espace,
        $d_{\mathrm{TV}}(\mu, \nu) \leq \P(X \neq Y)$.
  \item Calculer exactement, pour $p \in \intoo01$~:
        \[
        d_{\mathrm{TV}}\bigl(\mathcal B(1, p), \mathcal
        P(p)\bigr) = p\bigl(1 - \eu^{-p}\bigr) \leq p^2 .
        \]
  \item (Le Cam, par échange) Soient $X_i \sim \mathcal
        B(1, p_i)$ et $Y_i \sim \mathcal P(p_i)$, les $2n$
        variables \omterm{def:b3:probability:independence}{indépendantes}~; $S = X_1 + \dots + X_n$,
        et rappeler $Y_1 + \dots + Y_n \sim \mathcal
        P(\lambda)$ avec $\lambda = \sum_ip_i$
        (\cref{exo:b3:clt:1}). Échanger une coordonnée à la
        fois dans les hybrides entiers $H_i = Y_1 + \dots +
        Y_i + X_{i+1} + \dots + X_n$~: montrer, pour tout
        $A \subseteq \N$,
        \[
        \abs{\P(H_{i-1} \in A) - \P(H_i \in A)} \leq
        d_{\mathrm{TV}}\bigl(\mathcal B(1, p_i), \mathcal
        P(p_i)\bigr),
        \]
        et conclure l'\emph{inégalité de Le Cam}~:
        \[
        d_{\mathrm{TV}}\bigl(\text{loi de } S,\ \mathcal
        P(\lambda)\bigr) \leq \sum_{i=1}^np_i^2 .
        \]
  \item Dividendes. (a) Pour $p_i = \frac\lambda n$~: la
        borne est $\frac{\lambda^2}n$ --- la \omterm{def:b3:probability:space}{loi} des
        événements rares (\cref{exo:b3:clt:5}) montée à un
        taux explicite, uniforme sur tous les événements, et
        valable pour des $p_i$ inégaux aussi. (b) $500$
        lettres sont livrées, chacune s'égarant
        indépendamment avec probabilité $\frac1{500}$~:
        borner l'erreur du modèle de Poisson de paramètre
        $1$, et estimer la probabilité qu'aucune lettre ne
        s'égare. (c) Clore le problème~: comparer les deux
        classes d'universalité rencontrées ici ---
        gaussienne (beaucoup de petites contributions
        étalées~; deux moments appariés~; Taylor) et
        Poisson (beaucoup de contributions rares~; une
        moyenne appariée~; un couplage exact en variation
        totale) --- et la seule méthode de remplacement
        derrière les deux.

  \item (Erreur relative et transformée log) Soient
        $(X_n)$ i.i.d., positives, d'\omterm{def:b3:probability:space}{espérance} $\mu > 0$,
        de variance $\sigma^2$, et $\bar X_n$ la moyenne
        empirique. Montrer par la méthode delta que
        \[
        \sqrt n\,\bigl(\ln\bar X_n - \ln\mu\bigr)
        \Longrightarrow
        \mathcal N\Bigl(0,\ \frac{\sigma^2}{\mu^2}\Bigr) :
        \]
        le paramètre asymptotique de $\ln\bar X_n$ est le
        \emph{coefficient de variation} $\sigma/\mu$ ---
        erreur relative, sans échelle. En déduire un
        intervalle de confiance $95\%$ pour $\mu$ de la
        forme multiplicative $\bar
        X_n\cdot\eu^{\pm1.96\,\sigma/(\mu\sqrt n)}$, et
        expliquer quand il est préférable à l'additif.
  \item (Le troisième moment oriente l'erreur) Pour $X
        \sim$ Bernoulli($p$) centrée, calculer
        $\E\bigl[(X - p)^3\bigr] = p(1-p)(1-2p)$. En
        utilisant l'analyse de la Partie~IV (l'erreur
        d'échange est pilotée par les moments d'ordre
        trois), expliquer pourquoi l'approximation normale
        de $\mathcal B(n, p)$ est asymétrique pour $p \neq
        \frac12$ --- dépassant d'un côté, sous-estimant de
        l'autre --- et pourquoi $p = \frac12$ bénéficie du
        taux plus rapide à moments appariés. Vérifier le
        signe de l'asymétrie numériquement sur $\mathcal
        B(20, 0.1)$ contre $\mathcal N(2, 1.8)$~: comparer
        $\P(S = 0) = 0.9^{20}$ avec la masse gaussienne de
        $\intoo{-\infty}{0.5}$.
\end{enumerate}
\end{problem}
